\chapter{Sargantana Pipeline Overview}

\section{Pipeline Architecture}

The processor uses register renaming, speculative execution, branch prediction, instruction
queues, multiple execution units, and an in-order commit mechanism. These features allow
instructions to execute out of program order while preserving the correct architectural state at
retirement.

The complete processor pipeline is illustrated in Figure~\ref{fig:pipeline-overview}.

\begin{figure}[htbp]
  \centering
  \includegraphics[width=\textwidth]{ch02_pipeline_overview.png}
  \caption{Sargantana complete processor pipeline}
  \label{fig:pipeline-overview}
\end{figure}

The major stages of the pipeline are:

\begin{itemize}
  \item Fetch 1
  \item Fetch 2
  \item Decode
  \item Rename
  \item Register Read
  \item Execution
  \item Write Back
  \item Commit
\end{itemize}

Each stage performs a specific task and communicates with adjacent stages through dedicated
pipeline registers.

\section{Pipeline Stages}

The Sargantana pipeline consists of eight primary stages, each responsible for a particular
phase of instruction execution.

\subsection*{Fetch 1}
The Fetch 1 stage maintains the Program Counter (PC) and communicates with the instruction
cache to request instructions. Branch prediction hardware, including the Branch Target Buffer
(BTB) and Branch History Table (BHT), is also consulted to predict future control flow before
branch instructions are resolved.

\subsection*{Fetch 2}
The second fetch stage receives instructions from the instruction cache and forwards them to
the decode stage through the F2-to-D pipeline register. During this stage, prediction information
generated in Fetch 1 accompanies the fetched instructions.

\subsection*{Decode}
The decode stage interprets the binary instruction encoding and determines:
\begin{itemize}
  \item instruction type
  \item source registers
  \item destination register
  \item immediate values
  \item execution unit selection
  \item control signals
\end{itemize}
The decoded instruction is then inserted into the instruction queue for subsequent processing.

\subsection*{Rename}
One of the defining characteristics of Sargantana is its register renaming stage. Architectural
registers are mapped onto physical registers to eliminate false dependencies such as
Write-After-Write (WAW) and Write-After-Read (WAR) hazards.

This stage also interacts with:
\begin{itemize}
  \item Renaming Tables
  \item Free Lists
  \item Recovery Logic
\end{itemize}
to manage physical register allocation and pipeline recovery following branch mispredictions or
exceptions.

\subsection*{Register Read}
After renaming, source operands are read from the physical register file. If the required
operands have recently been produced but not yet written back, bypass logic forwards the
values directly from later pipeline stages to reduce execution latency.

\subsection*{Execution}
Instructions are dispatched to specialized execution units depending on their operation.
These include:
\begin{itemize}
  \item Arithmetic Logic Unit (ALU)
  \item Branch Unit (BrU)
  \item Load/Store Unit (Memory)
  \item Multiply/Divide Unit
  \item Floating Point Unit (FPU)
  \item SIMD Unit
\end{itemize}
Each functional unit operates independently, allowing several instructions to execute
concurrently.

\subsection*{Write Back}
Upon completion of execution, results are written back into the physical register file. These
results become available for dependent instructions through both register reads and forwarding
paths.

\subsection*{Commit}
Although instructions execute out-of-order, they retire strictly in program order. The commit
stage updates the architectural state only after ensuring that all previous instructions have
completed successfully. This guarantees precise exceptions and preserves program
correctness.

\section{Instruction Flow through the Pipeline}

Figure~\ref{fig:pipeline-overview} illustrates how instructions travel through the processor pipeline.

Unlike an in-order processor, instructions may execute as soon as their operands become
available rather than waiting for older instructions to finish execution. However, all instructions
commit in the original program order, ensuring correct architectural behavior.
